\subsection{Asociaciones de recursos con actividades existentes}\label{sec:impl-asociaciones-recursos}

La asociación organizativa permite relacionar un recurso declarado, y opcionalmente uno de sus actos, con una audiencia o un plazo que ya existen en el mismo expediente. El vínculo tiene identidad e historia independientes de sus extremos. Su registro no programa otra audiencia, recalcula un término ni genera una segunda alerta. La decisión se documenta en \rutarepo{docs/adr/0041-exact-resource-activity-associations.md} y el intercambio HTTP en \rutarepo{docs/resource-activities-api.md}.

\subsubsection{Selección exacta y cabeza esperada}

La selección conserva identidad, revisión y digest del recurso. Cuando incluye un acto, identifica también la revisión del recurso que contiene esa captura. Ambas selecciones son independientes: puede vincularse la revisión 1 del recurso junto con un acto conservado en la revisión 2, aunque una corrección posterior sea la cabeza actual. El servidor verifica que cada fuente pertenezca al expediente y al recurso solicitados; no sustituye la selección por la última revisión ni vuelve a admitir el contenido documental retenido.

La audiencia ordinaria se identifica por su revisión y digest de envío; el plazo y la audiencia propia de un recurso, por su revisión y digest de captura. La cabeza esperada del recurso se declara aparte de esas referencias históricas. Preparar reúne las fuentes verificadas y el contexto administrativo; confirmar vuelve a contrastarlos bajo el bloqueo compartido de auditoría. Una revisión concurrente produce conflicto y exige aceptar expresamente otra base.

La vinculación crea la revisión 1. La desvinculación añade una revisión 2 con motivo, preservando selección y fuentes. Es terminal para esa identidad; un vínculo posterior tiene otra raíz. Pueden existir varias asociaciones con una misma actividad sin duplicar la actividad ni sus avisos. El archivo del recurso impide crear vínculos nuevos, pero permite desvincular los existentes cuando el expediente está activo.

\subsubsection{Historia y estado actual separados}

Detalle y lista exponen dos partes: la asociación histórica con sus capturas, y la cabeza actual verificada de la actividad. Ambas se resuelven en una transacción autorizada. Una audiencia histórica puede conservar su programación original aunque la cabeza actual esté cancelada. Para un plazo, la evaluación capturada permanece legible, pero sólo la proyección operativa de la cabeza actual determina si se muestra un vencimiento vigente. Un plazo retirado, pendiente o afectado por cambios no recupera una fecha operativa desde su cálculo histórico.

La consulta toma un instante UTC común después de adquirir el bloqueo. Todas las filas de la página y las proyecciones operativas de plazo comparten ese corte. El servicio comprueba que esté entre el inicio de la solicitud y el retorno del almacén. Así admite una cabeza confirmada mientras el lector esperaba, sin aceptar una observación anterior a la solicitud o posterior al retorno. La historia y los recibos de escritura no incorporan esa afirmación de vigencia.

\subsubsection{Persistencia, permisos e interfaz}

Las migraciones de asociaciones, conservadas en el directorio \rutarepo{migrations/}, incorporan dos tablas: raíces y revisiones de asociación. La transacción añade revisión, recibo y auditoría conjuntamente. El catálogo estricto comprueba estructura y privilegios; el inventario reconstruye capturas históricas en lotes acotados, sin modificar datos al abrir. El respaldo conserva ambas tablas y sus dependencias, no sólo los identificadores de los extremos. La \cref{sec:pruebas-asociaciones-recursos} reporta la aceptación API con restauración, que comparó capturas y respuestas y validó por separado los instantes de consulta nuevos, junto con los recorridos de navegador real.

Owner gestiona todos los expedientes; Litigante asignado gestiona y Paralegal asignado consulta. Cliente no accede. El servicio reautentica al principal completo y el adaptador comprueba cuenta, rol y asignación vigentes. El cierre conserva lecturas y repetición exacta autorizada de una operación confirmada, pero impide nuevas mutaciones.

Qadra presenta selectores de actividad, revisión y acto opcional, revisión previa a confirmar, listado filtrado, detalle e historia; la captura vinculada y la actividad actual ocupan áreas separadas. Conflictos, respuestas inciertas y revocación siguen los mecanismos comunes de las \cref{sec:impl-qadra-borradores,sec:impl-qadra-conciliacion}: un conflicto conserva selección y motivo hasta aceptar una nueva base, ninguna coincidencia visual sustituye al recibo exacto, y el reenvío explícito tras un resultado incierto conserva el mismo comando, operación y digest, sin preparar otra identidad.

Las asociaciones sólo vinculan actividades que ya existen. La creación conjunta de un plazo y su asociación tiene el contrato independiente de la \cref{sec:impl-plazo-contextual}, y la programación de una audiencia propia crea su asociación inicial en la misma transacción (\cref{sec:impl-audiencias-recursos}). La consulta inversa de la \cref{sec:impl-recursos-desde-actividad} completa un recorrido de lectura desde actividades y avisos, sin crear actividades adicionales. Los vínculos no activan términos ni amplían la cobertura jurídica delimitada al final de la \cref{sec:impl-recursos-declarados}, y no modifican los objetivos aprobados.

\subsubsection{Consulta inversa y regreso a la actividad}\label{sec:impl-recursos-desde-actividad}

La consulta inversa parte de una audiencia ordinaria o de un plazo identificado dentro del expediente y devuelve las asociaciones actuales de sus recursos. Utiliza las rutas de audiencia y plazo con sufijo \texttt{resource-associations}, sin recorrer en el navegador todos los recursos del expediente. La lista muestra vínculos activos de manera predeterminada, permite consultar los desvinculados o ambos estados y pagina por UUID ascendente con un límite de 1 a 100. El estado se filtra después de seleccionar la cabeza de cada asociación y antes de aplicar el límite. Cada continuación conserva expediente, destino y filtro; las páginas son observaciones independientes, no una instantánea histórica retenida entre peticiones.

Cada respuesta incluye un instante UTC común incluso cuando no contiene filas. La transacción comprueba el destino, las capturas de recurso, acto y actividad, sus cabezas actuales y la auditoría de lectura; el servicio reautentica al principal antes de entregar datos. Owner consulta todos los expedientes y Litigante o Paralegal necesitan asignación vigente. Cliente permanece denegado. El cierre del expediente y el archivo del recurso conservan las lecturas autorizadas. Un error de acceso no se representa como ausencia de relaciones.

Qadra presenta el panel \enquote{Recursos relacionados} en el detalle de actividad. Se distinguen la revisión abierta, la revisión de actividad capturada por cada vínculo y la cabeza actual observada. La lista representa el estado actual de las asociaciones: una relación registrada después de emitir una alerta puede aparecer, sin afirmar que existiera entonces ni que causara el aviso. La selección conserva la identidad, revisión y digest tanto del recurso como de la asociación, y su apertura comprueba las capturas exactas. Consultar el recurso actual sigue siendo una acción diferente.

El recorrido conserva el regreso a la revisión de actividad desde la que se abrió el recurso. Si comenzó en una alerta, mantiene también el regreso a la bandeja y sus filtros. Las consultas no marcan el aviso como leído ni modifican su origen; la revocación y las respuestas tardías siguen los mecanismos comunes de las \cref{sec:impl-qadra-borradores,sec:impl-qadra-conciliacion}. El contrato está en \rutarepo{docs/activity-resource-links-api.md} y la decisión en \rutarepo{docs/adr/0059-activity-resource-navigation.md}. Su verificación se distingue de la evidencia de las asociaciones en la \cref{sec:pruebas-recursos-desde-actividad}.
